\begin{frame}{Ejemplo Numérico del Procedimiento de Cálculo}

    \begin{columns}[T]
        %========================================
        % COLUMNA IZQUIERDA: LOS PASOS LÓGICOS
        %========================================
        \column{0.46\textwidth}
        \begin{block}{Mecánica de Evaluación}
            \scriptsize
            Para cada punto candidato $(x,y,z)$, el algoritmo ejecuta tres pasos secuenciales por cada sensor:
            \begin{enumerate}
                \item \textbf{Geometría:} Calcula la distancia euclidiana 3D al sensor.
                \item \textbf{Física:} Evalúa la atenuación teórica ($A_{pred}$).
                \item \textbf{Estadística:} Obtiene el residuo individual.
            \end{enumerate}
        \end{block}

        \vspace{0.1cm}

        \begin{block}{Acumulación del Costo}
            \scriptsize
            Este procedimiento matemático se repite de forma idéntica para los 12 sensores de la red. La sumatoria de estos residuos cuadrados es la que define el error global $E(x,y,z)$.
        \end{block}

        %========================================
        % COLUMNA DERECHA: LA CORRIDA EN FRÍO REAL (S9)
        %========================================
        \column{0.50\textwidth}
        \begin{exampleblock}{Traza de Ejecución Real (Sensor 9)}
            \tiny
            Evaluación en el punto estimado: $(1.2001, -0.7765, -1.1381)$ \\
            Datos de $S_9$: Coordenadas $(0.5, 0.5, 0.0)$ y $Az_{obs} = 0.852201$
            
            \vspace{0.15cm}
            \textbf{Paso 1: Distancia Euclidiana ($R_9$)}
            \[ R_9 = \sqrt{(0.5 - 1.2001)^2 + (0.5 - (-0.7765))^2 + (0.0 - (-1.1381))^2} \]
            \[ R_9 = \sqrt{(-0.7001)^2 + (1.2765)^2 + (1.1381)^2} \approx 1.8479 \]

            \vspace{0.1cm}
            \textbf{Paso 2: Amplitud Predicha ($A_{pred,9}$)}
            \[ A_{pred,9} = 10 \cdot \frac{e^{-1.8479}}{1.8479} \approx 0.852651 \]

            \vspace{0.1cm}
            \textbf{Paso 3: Residuo Individual}
            \[ \text{Residuo} = A_{obs,9} - A_{pred,9} \]
            \[ \text{Residuo} = 0.852201 - 0.852651 = -0.000450 \]
        \end{exampleblock}
    \end{columns}

\end{frame}